% GENERATED FILE. Edit canonical lessons, then run npm run generate:books.
% canonical-source-hash: fnv1a64:0f4653b94c6efb30
% canonical-lessons: GE-C374-entschliessen, GE-C374-entschluss, GE-C374-ausweg, GE-C374-flexibel, GE-C374-planung

\chapter{Sich entschließen}
\label{ch:ge-sich-entschlie-en}
\hlchaptermodality{374}

\begin{chapteropening}
\textbf{By the end of this chapter:} \emph{I can say that I have made up my mind, find a way out and plan in a flexible way.}

Complete the last lesson of chapter 374: I can say that I have made up my mind, find a way out and plan in a flexible way.
\end{chapteropening}

\section[entschließen]{entschließen — to make up your mind (reflexive, sich entschließen)}

\label{lesson:GE-C374-entschliessen}



\begin{quote}
Before the new one: say the German for the forecast, then the German for to allow for.
\end{quote}

\subsection*{You'll want to know: entschließen}
\textbf{entschließen} — \textquotedblleft{}to make up your mind (reflexive, sich entschließen)\textquotedblright{}.

\begin{quote}
Nach langem Nachdenken hat sie sich entschlossen, in eine größere Stadt zu ziehen.
\end{quote}

\begin{quote}
\emph{After thinking it over for a long time, she made up her mind to move to a bigger city.}
\end{quote}

Reflexive and strong: \emph{sie entschließt sich}, \emph{sie hat sich entschlossen}. \emph{Entschlossen} — determined.

\begin{cousinweb}
\emph{entschließen} is built with \emph{ent-} on \emph{schließen} \textquotedblleft{}to close, to conclude\textquotedblright{}, inherited from Old High German \emph{sliozan} and Proto-Germanic \emph{sleutaną} (EtymDB): you close a question.
\end{cousinweb}

\begin{grammarlens}[title={What you've built}]
One more word for telling, explaining and describing at B1.
\end{grammarlens}

\subsection*{Your turn}
\begin{itemize}
\raggedright
  \item \emph{Say it:} \emph{entschließen}
  \item \emph{Say it:} \emph{entschließen}, once more
  \item \emph{Say it:} a sentence of your own with \emph{entschließen}
  \item \emph{Recall:} say the German for the forecast, then the German for to allow for, then say \emph{entschließen} again
\end{itemize}

\begin{tcolorbox}[breakable,colback=teal!4,colframe=teal!35!black,title={Before you move on}]
What does entschließen mean? (\textquotedblleft{}To make up your mind (reflexive, sich entschließen)\textquotedblright{}.) Say it once more.
\end{tcolorbox}
\section[der Entschluss]{der Entschluss — the resolution, the decision}

\label{lesson:GE-C374-entschluss}



\begin{quote}
Before the new one: say the German for to allow for, then the German for to make up your mind.
\end{quote}

\subsection*{You'll want to know: der Entschluss}
\textbf{der Entschluss} — \textquotedblleft{}the resolution, the decision\textquotedblright{}.

\begin{quote}
Mein Entschluss steht fest: Ich höre auf zu rauchen.
\end{quote}

\begin{quote}
\emph{My mind is made up: I am giving up smoking.}
\end{quote}

Plural: \emph{die Entschlüsse}. \emph{Einen Entschluss fassen} — to make a decision.

\begin{cousinweb}
\emph{Entschluss} is formed from \emph{entschließen} (EtymDB), built on \emph{schließen} \textquotedblleft{}to close\textquotedblright{}, Old High German \emph{sliozan}, Proto-Germanic \emph{sleutaną} (EtymDB).
\end{cousinweb}

\begin{grammarlens}[title={What you've built}]
One more word for telling, explaining and describing at B1.
\end{grammarlens}

\subsection*{Your turn}
\begin{itemize}
\raggedright
  \item \emph{Say it:} \emph{der Entschluss}
  \item \emph{Say it:} \emph{der Entschluss}, once more
  \item \emph{Say it:} a sentence of your own with \emph{der Entschluss}
  \item \emph{Recall:} say the German for to allow for, then the German for to make up your mind, then say \emph{der Entschluss} again
\end{itemize}

\begin{tcolorbox}[breakable,colback=teal!4,colframe=teal!35!black,title={Before you move on}]
What does der Entschluss mean? (\textquotedblleft{}The resolution, the decision\textquotedblright{}.) Say it once more.
\end{tcolorbox}
\section[der Ausweg]{der Ausweg — the way out}

\label{lesson:GE-C374-ausweg}



\begin{quote}
Before the new one: say the German for to make up your mind, then the German for the resolution.
\end{quote}

\subsection*{You'll want to know: der Ausweg}
\textbf{der Ausweg} — \textquotedblleft{}the way out\textquotedblright{}.

\begin{quote}
Es gibt immer einen Ausweg.
\end{quote}

\begin{quote}
\emph{There is always a way out.}
\end{quote}

Plural: \emph{die Auswege}. \emph{Der letzte Ausweg} — the last resort.

\begin{cousinweb}
\emph{Ausweg} joins \emph{aus} \textquotedblleft{}out\textquotedblright{} to \emph{Weg} \textquotedblleft{}way\textquotedblright{}, inherited from Old High German \emph{weg} and Proto-Germanic \emph{wegaz} \textquotedblleft{}road, path\textquotedblright{} (EtymDB).
\end{cousinweb}

\begin{grammarlens}[title={What you've built}]
One more word for telling, explaining and describing at B1.
\end{grammarlens}

\subsection*{Your turn}
\begin{itemize}
\raggedright
  \item \emph{Say it:} \emph{der Ausweg}
  \item \emph{Say it:} \emph{der Ausweg}, once more
  \item \emph{Say it:} a sentence of your own with \emph{der Ausweg}
  \item \emph{Recall:} say the German for to make up your mind, then the German for the resolution, then say \emph{der Ausweg} again
\end{itemize}

\begin{tcolorbox}[breakable,colback=teal!4,colframe=teal!35!black,title={Before you move on}]
What does der Ausweg mean? (\textquotedblleft{}The way out\textquotedblright{}.) Say it once more.
\end{tcolorbox}
\section[flexibel]{flexibel — flexible}

\label{lesson:GE-C374-flexibel}



\begin{quote}
Before the new one: say the German for the resolution, then the German for the way out.
\end{quote}

\subsection*{You'll want to know: flexibel}
\textbf{flexibel} — \textquotedblleft{}flexible\textquotedblright{}.

\begin{quote}
Bei der Arbeitszeit bin ich ganz flexibel.
\end{quote}

\begin{quote}
\emph{I am very flexible about working hours.}
\end{quote}

With an ending the \emph{e} drops out: \emph{flexible Arbeitszeiten}.

\begin{cousinweb}
\emph{flexibel} is a learned borrowing from Latin \emph{flexibilis} \textquotedblleft{}that may be bent, pliant\textquotedblright{} (EtymDB; Lewis \& Short), from \emph{flectere} \textquotedblleft{}to bend\textquotedblright{}.
\end{cousinweb}

\begin{grammarlens}[title={What you've built}]
One more word for telling, explaining and describing at B1.
\end{grammarlens}

\subsection*{Your turn}
\begin{itemize}
\raggedright
  \item \emph{Say it:} \emph{flexibel}
  \item \emph{Say it:} \emph{flexibel}, once more
  \item \emph{Say it:} a sentence of your own with \emph{flexibel}
  \item \emph{Recall:} say the German for the resolution, then the German for the way out, then say \emph{flexibel} again
\end{itemize}

\begin{tcolorbox}[breakable,colback=teal!4,colframe=teal!35!black,title={Before you move on}]
What does flexibel mean? (\textquotedblleft{}Flexible\textquotedblright{}.) Say it once more.
\end{tcolorbox}
\section[die Planung]{die Planung — the planning}

\label{lesson:GE-C374-planung}



\begin{quote}
Before the new one: say the German for the way out, then the German for flexible.
\end{quote}

\subsection*{You'll want to know: die Planung}
\textbf{die Planung} — \textquotedblleft{}the planning\textquotedblright{}.

\begin{quote}
Die Planung der Hochzeit hat ein Jahr gedauert.
\end{quote}

\begin{quote}
\emph{Planning the wedding took a year.}
\end{quote}

Plural: \emph{die Planungen}. \emph{In Planung} — at the planning stage.

\begin{cousinweb}
\emph{Planung} is the noun of \emph{planen}, from \emph{Plan}, borrowed from French \emph{plan} and going back to Latin \emph{plānus} \textquotedblleft{}level, flat\textquotedblright{} (EtymDB).
\end{cousinweb}

\begin{grammarlens}[title={What you've built}]
One more word for telling, explaining and describing at B1.
\end{grammarlens}

\subsection*{Your turn}
\begin{itemize}
\raggedright
  \item \emph{Say it:} \emph{die Planung}
  \item \emph{Say it:} \emph{die Planung}, once more
  \item \emph{Say it:} a sentence of your own with \emph{die Planung}
  \item \emph{Recall:} say the German for the way out, then the German for flexible, then say \emph{die Planung} again
\end{itemize}

\begin{tcolorbox}[breakable,colback=teal!4,colframe=teal!35!black,title={Before you move on}]
What does die Planung mean? (\textquotedblleft{}The planning\textquotedblright{}.) Say it once more.
\end{tcolorbox}
